\documentclass[11pt]{article}
\usepackage[margin=1in]{geometry}
\usepackage{amsmath,amssymb,amsthm}
\usepackage{hyperref}

\newtheorem{theorem}{Theorem}
\newtheorem{lemma}{Lemma}
\newtheorem{corollary}{Corollary}
\newtheorem{definition}{Definition}
\newtheorem{remark}{Remark}
\newtheorem{example}{Example}

% ---- notation shortcuts (kept consistent with three_theorems_proofs.tex) ----
\newcommand{\adj}{\operatorname{adj}}
\newcommand{\spann}{\operatorname{span}}
\newcommand{\Ran}{\operatorname{Ran}}
\newcommand{\kry}{\mathcal K}
\newcommand{\nobs}{\mathcal N_{\mathrm{obs}}}
\newcommand{\Vresp}{\mathcal V_{\mathrm{resp}}}
\newcommand{\nudiss}{\nu_{\mathrm{diss}}}

% ---- a simple "take-away" box that needs no extra packages ----
\newenvironment{takeaway}
  {\par\medskip\noindent
   \begin{center}\begin{minipage}{0.92\textwidth}
   \hrule\medskip\noindent\textbf{Take-away.}\ }
  {\medskip\hrule\end{minipage}\end{center}\medskip}

\title{The New No-Go Theory, Gently:\\
An Undergraduate's Guide to the Three-Class\\
Classification of Sector-Mediated Response}
\author{}
\date{}

\begin{document}
\maketitle

\begin{abstract}
This is a pedagogical companion to the research documents
\emph{Three Central Theorems of the Sector-Resolved Response Classification}
and \emph{A Closed Theorem Package for Master-Response No-Go Classification}.
Those documents prove, at full rigor, that the question ``does a designated
coherent pathway of an open quantum system contribute to what a detector
sees?'' has exactly three possible answers --- \emph{never} (an exact
structural zero), \emph{less and less} (algebraic suppression under strong
dissipation), or \emph{always} (a protected response that survives arbitrarily
strong dissipation) --- and that which answer applies can be decided by a
finite linear-algebra computation.  This guide explains the same theory using
only tools from a standard undergraduate curriculum: matrix inverses,
eigenvalues, geometric series, and block matrices.  Every technical
simplification made here is flagged explicitly; nothing is silently
misstated.  After reading this guide you should be able to open the rigorous
documents and recognize every object in them.
\end{abstract}

\tableofcontents

% ==================================================================
\section{What is a no-go theorem, and why do we need a new one?}
% ==================================================================

\subsection{No-go theorems in general}

A \emph{no-go theorem} is a mathematical proof that something is impossible
within a stated set of rules.  Famous examples: the \emph{no-cloning theorem}
(no quantum machine can copy an unknown quantum state) and
\emph{Earnshaw's theorem} (no arrangement of static charges can trap a charge
stably).  A good no-go theorem does two things at once: it saves you from
wasting effort on impossible designs, and --- by listing its assumptions
precisely --- it tells you exactly which assumption you must break if you want
the ``impossible'' effect after all.

The theory explained here concerns a specific physical question from quantum
optics.  In effects like \emph{electromagnetically induced transparency}
(EIT), a weak probe laser interacts with atoms (or solid-state defects such as
color centers), and quantum interference between internal pathways can make an
opaque medium transparent.  A natural engineering question is:

\begin{center}
\emph{``Does this particular internal pathway (`sector') of my system\\
actually contribute to the signal my detector measures?''}
\end{center}

An older result --- referred to below as \emph{the old no-go criterion} ---
answered this with a binary yes/no.  The new theory shows that the binary
answer is too coarse: the honest answer has \textbf{three} tiers, and the old
criterion cannot distinguish two of them.

\subsection{The three tiers, in plain words}

Imagine you can turn a single knob $\Gamma$ that makes the environment noisier
and noisier --- more dissipation, more decoherence, everything decaying
faster.  Watch what happens to the part of the detector signal that is caused
by the designated pathway.  The new theory proves that exactly one of three
things happens:

\begin{description}
\item[Class I (exact structural no-go).]  The pathway contributes
\emph{exactly zero}, at \emph{every} probe frequency, for \emph{every} value
of $\Gamma$.  Not small --- zero.  The wire is simply not connected.
\item[Class II (asymptotic suppression).]  The contribution is nonzero but
dies algebraically as the noise grows: it shrinks like $\Gamma^{-\nu}$ for
some positive integer $\nu$.  The signal must pass through $\nu$ heavily
damped ``rooms,'' and each door costs a factor $1/\Gamma$.
\item[Class III (protected survival).]  The contribution approaches a
\emph{finite nonzero limit} as $\Gamma\to\infty$.  There is a soundproof room
--- a \emph{dark subspace} the dissipation cannot touch --- and the
conversation inside continues at full volume no matter how loud the corridor
gets.
\end{description}

The whole classification is summarized by a single number, the
\emph{response-class index}
\[
\nu\;\in\;\{\infty\}\;\cup\;(0,\infty)\;\cup\;\{0\},
\]
where $\nu=\infty$ means Class I, $0<\nu<\infty$ means Class II
(with $\nu$ the decay exponent), and $\nu=0$ means Class III.

\begin{takeaway}
The new no-go theory upgrades a yes/no verdict to a three-way classification
--- exact zero, power-law suppressed, or protected --- indexed by one number
$\nu$, and proves the three cases are mutually exclusive, exhaust all
possibilities (within the stated scope), and can be decided by a finite
computation.
\end{takeaway}

\subsection{What you need to know already}

Only standard undergraduate material:
\begin{itemize}
\item matrices, matrix inverses, eigenvalues and eigenvectors;
\item the geometric series $1/(1-x)=\sum_k x^k$ for $|x|<1$;
\item block (partitioned) matrices;
\item the words ``analytic function'' and ``Laurent series'' (we use them
only descriptively);
\item it helps --- but is not required --- to have seen a driven damped
harmonic oscillator.
\end{itemize}
You do \emph{not} need to know open-quantum-system theory.  Wherever a
concept from that field appears (Lindblad equation, Liouville space, dark
state), we describe it in words and then work purely with matrices.

% ==================================================================
\section{The playground: poke the system, listen to the response}
% ==================================================================

\subsection{Warm-up: one damped oscillator}

Drive a damped harmonic oscillator with a force $F e^{-i\omega t}$ and measure
its steady-state displacement.  The answer has the form
\[
x(\omega) \;=\; \underbrace{\frac{1}{\;\omega_0^2-\omega^2-i\gamma\omega\;}}_{\text{transfer function}}\,F .
\]
Everything interesting --- resonances, damping, interference --- lives in the
\emph{transfer function}: the complex number that converts ``how hard you
poke'' into ``how much you hear,'' as a function of frequency.  All of linear
response theory is a matrix generalization of this single formula.

\subsection{The general setup: a matrix, a source, a readout}

An open quantum system probed by a weak laser can be described (this is the
one imported fact from quantum optics; see the box below) by a family of
matrices
\begin{equation}
A_\Gamma(z) \;=\; \Gamma D \;+\; B(z),
\label{eq:AGamma}
\end{equation}
acting on a finite-dimensional complex vector space
$\mathcal V\cong\mathbb C^N$, where:
\begin{itemize}
\item $z$ is the complex frequency variable (think ``probe detuning'');
\item $B(z)$ collects the \emph{coherent} physics: energy levels, laser
couplings, detunings, and any dissipation whose rate is held fixed.  It
depends analytically on $z$;
\item $D$ is the \emph{dissipation matrix}: the part of the dynamics whose
overall strength we will crank up.  It does not depend on $z$;
\item $\Gamma\ge 0$ is the knob: the overall dissipation scale.
\end{itemize}
Two fixed vectors complete the cast:
\begin{itemize}
\item $c\in\mathbb C^N$, the \emph{source}: where the probe field injects
excitation (``where you poke'');
\item $p\in\mathbb C^N$, the \emph{readout}: which coherence the detector
measures (``where you listen'').
\end{itemize}
The measured linear response is the scalar transfer function
\begin{equation}
F_\Gamma(z) \;=\; p^\dagger\, A_\Gamma(z)^{-1}\, c .
\label{eq:FGamma}
\end{equation}
Read it from right to left: inject at $c$, let the system respond
(that is what the inverse matrix does --- it solves the steady-state
equations), and project the result onto what the detector sees.  For the
single oscillator above, $A$ is $1\times 1$ and
\eqref{eq:FGamma} reduces exactly to the warm-up formula.

\begin{takeaway}
One matrix family $A_\Gamma(z)=\Gamma D+B(z)$, one source vector $c$, one
readout vector $p$.  Everything the detector can see at linear order is the
scalar $F_\Gamma(z)=p^\dagger A_\Gamma(z)^{-1}c$.  The entire theory is about
how this scalar behaves as $\Gamma\to\infty$.
\end{takeaway}

\begin{remark}[Where equation \eqref{eq:AGamma} comes from --- optional]
For readers who have met density matrices: a Markovian open quantum system
evolves under a Lindblad master equation
$\dot\rho=\mathcal L\rho$.  The superoperator $\mathcal L$ is a linear map on
the space of matrices $\rho$; flattening $\rho$ into a long column vector
turns $\mathcal L$ into an ordinary $N\times N$ matrix ($N=d^2$ for a
$d$-level system).  This flattened arena is called \emph{Liouville space}.
Splitting $\mathcal L$ into the dissipators we scale ($\Gamma D$) and
everything else ($B$), and solving the weak-probe steady state at frequency
$z$, produces exactly \eqref{eq:AGamma}--\eqref{eq:FGamma}.  One important
consequence: $D$ acts on Liouville space and is \emph{in general not
Hermitian}, even though the physics is perfectly unitary-plus-dissipative.
This is why Section~\ref{sec:thm3} must use slightly more careful projections
than the orthogonal ones you know from inner-product spaces.
\end{remark}

\subsection{The sector cut and the master response}
\label{sec:master}

Now the key construction.  We want to know what a designated internal pathway
--- a \emph{sector} $S$, e.g.\ ``the coupling through the metastable spin
state'' --- contributes to the signal.  The clean way to ask this: compare the
real system with a fictitious surgery on it.

\begin{definition}[Sector cut]
The \emph{cut model} is obtained from the full model by deleting the
couplings that connect sector $S$ to the rest, \emph{inside $B$ only}: the
dissipation matrix $D$, the source $c$, and the readout $p$ are left
untouched.  Write $B_{\mathrm{full}}(z)$ and $B^{(S)}_{\mathrm{cut}}(z)$ for
the two coupling matrices.
\end{definition}

\begin{definition}[Master sector-resolved response]
The \emph{master response} of sector $S$ is the difference of the two
transfer functions,
\begin{equation}
R_{S,\Gamma}(z)\;=\;\chi_{\mathrm{full},\Gamma}(z)-\chi^{(S)}_{\mathrm{cut},\Gamma}(z)
\;=\;p^\dagger\Big[\big(\Gamma D+B_{\mathrm{full}}(z)\big)^{-1}
-\big(\Gamma D+B^{(S)}_{\mathrm{cut}}(z)\big)^{-1}\Big]c .
\label{eq:master}
\end{equation}
\end{definition}

$R_{S,\Gamma}(z)$ is precisely ``the part of the detected response that
disappears when you unplug sector $S$.''  If $R_{S,\Gamma}\equiv 0$, the
sector is irrelevant to the detector; if it survives strong dissipation, the
sector-mediated effect is robust.  \textbf{The object being classified is
this difference} --- not the full response and the cut response separately.
This matters: the full and cut responses could each be large while their
difference cancels to zero, or each could decay while the difference decays
faster.  Classifying the difference directly keeps all such cancellations
honest.

\begin{remark}[One scalar to rule them all: the doubling trick]
A difference of two transfer functions looks like a more complicated object
than a single transfer function --- but it is not.  Stack two copies of the
state space into $\hat{\mathcal V}=\mathcal V\oplus\mathcal V$, run the full
model in the first copy and the cut model in the second, feed the same source
into both, and read out with $\hat p=(p,\,-p)$, i.e.\ \emph{subtract} the two
outputs at the readout stage:
\[
\hat D=\begin{pmatrix}D&0\\0&D\end{pmatrix},\quad
\hat B_S(z)=\begin{pmatrix}B_{\mathrm{full}}(z)&0\\0&B^{(S)}_{\mathrm{cut}}(z)\end{pmatrix},\quad
\hat c=\begin{pmatrix}c\\ c\end{pmatrix},\quad
\hat p=\begin{pmatrix}p\\ -p\end{pmatrix}.
\]
Then $R_{S,\Gamma}(z)=\hat p^\dagger(\Gamma\hat D+\hat B_S(z))^{-1}\hat c$
--- a single scalar transfer function of exactly the form
\eqref{eq:FGamma}, on a space of dimension $2N$.  Consequently, every theorem
we prove about scalar transfer functions $F_\Gamma$ automatically applies to
the master response $R_{S,\Gamma}$.  In the rest of this guide we therefore
state everything for a generic $F_\Gamma(z)=p^\dagger(\Gamma D+B(z))^{-1}c$
and use small matrices in the examples; keep in mind that in applications
this $F_\Gamma$ \emph{is} the master response in its doubled realization.
\end{remark}

\subsection{The one question}

Fix a compact window $K$ of physical frequencies (the range your laser can
scan).  Ask:
\begin{center}
\emph{What happens to $F_\Gamma(z)$, uniformly for $z\in K$, as
$\Gamma\to\infty$?}
\end{center}
The three theorems of the next sections answer this question in the three
possible regimes, and Section~\ref{sec:trichotomy} proves the regimes are the
only ones and never overlap.

% ==================================================================
\section{The toolbox: three facts from linear algebra}
\label{sec:toolbox}
% ==================================================================

Three elementary facts power the entire theory.  We state them here once so
the theorems read smoothly.

\paragraph{Tool 1: the Neumann series (geometric series for matrices).}
If $\|Y\|<1$ then $I+Y$ is invertible and
\[
(I+Y)^{-1}=I-Y+Y^2-Y^3+\cdots,
\]
with the same proof as for numbers (multiply and telescope; the series
converges absolutely because $\|Y^k\|\le\|Y\|^k$).  We will apply this with
$Y$ proportional to $1/\Gamma$: strong dissipation makes the correction term
small, so we can expand in powers of $1/\Gamma$.

\paragraph{Tool 2: Cayley--Hamilton.}
Every $n\times n$ matrix $M$ satisfies its own characteristic polynomial:
$M^n$ is a linear combination of $I,M,\dots,M^{n-1}$.  Consequence: the
infinite list of numbers $\mu_k=\ell^\dagger M^k r$ ($k=0,1,2,\dots$) is
completely determined by its first $n$ entries.  If
$\mu_0=\dots=\mu_{n-1}=0$, then \emph{all} $\mu_k$ vanish.  This is what
makes ``check infinitely many frequencies'' collapse to ``check $n$
numbers'' --- the origin of every \emph{finiteness} claim in the theory.

\paragraph{Tool 3: block elimination (the Schur complement).}
To solve
$\left(\begin{smallmatrix}A_{11}&A_{12}\\A_{21}&A_{22}\end{smallmatrix}\right)
\left(\begin{smallmatrix}x_1\\x_2\end{smallmatrix}\right)
=\left(\begin{smallmatrix}b_1\\b_2\end{smallmatrix}\right)$
when $A_{22}$ is invertible, solve the second row for $x_2$ and substitute
into the first:
\[
\big(\underbrace{A_{11}-A_{12}A_{22}^{-1}A_{21}}_{\text{Schur complement}}\big)\,x_1
= b_1-A_{12}A_{22}^{-1}b_2 .
\]
Physically: eliminate the fast, damped variables and obtain an effective
equation for the slow, protected ones.  This is the engine of
Theorem~III.

% ==================================================================
\section{Theorem I: the exact structural no-go (Class I, $\nu=\infty$)}
\label{sec:thm1}
% ==================================================================

\subsection{Setting: moments of a realization}

Any scalar transfer function of our type can be written (possibly after the
doubling trick and a change of variables) in the standard state-space form
\begin{equation}
\delta\chi_S(z)=\ell^\dagger (zI-M)^{-1} r,
\label{eq:realization}
\end{equation}
with an $n\times n$ matrix $M$, input vector $r$, and output vector $\ell$.
For $|z|>\|M\|$ the geometric series (Tool 1, with $Y=M/z$) gives
\[
\delta\chi_S(z)=\sum_{k\ge0}\frac{\mu_k}{z^{k+1}},
\qquad
\mu_k:=\ell^\dagger M^k r .
\]
The numbers $\mu_k$ are called \emph{moments}.  They have a beautiful
mechanical reading: $r$ is where the excitation enters; $M^k r$ is where it
has spread after $k$ steps of internal dynamics; $\mu_k$ asks whether, after
$k$ steps, any of it overlaps what the detector measures.

Two subspaces make this picture precise:
\begin{itemize}
\item the \emph{reachable subspace}
$\kry(M,r)=\spann\{r,\,Mr,\,\dots,\,M^{n-1}r\}$
--- everywhere the input can spread;
\item the \emph{unobservable subspace}
$\nobs(M,\ell)=\bigcap_{k\ge0}\ker(\ell^\dagger M^k)$
--- everything the detector can never see, no matter how long it waits.
\end{itemize}
These are the classical \emph{Kalman} subspaces from control theory.

\subsection{The theorem}

\begin{theorem}[Exact structural no-go]
\label{thm:one}
For the realization \eqref{eq:realization} the following are equivalent:
\begin{enumerate}
\item[(i)] $\mu_k=\ell^\dagger M^k r=0$ for $k=0,1,\dots,n-1$
\hfill (\emph{$n$ numbers vanish});
\item[(ii)] $\kry(M,r)\subseteq\nobs(M,\ell)$
\hfill (\emph{everything reachable is invisible});
\item[(iii)] $\ell^\dagger q(M)\,r=0$ for every polynomial $q$
\hfill (\emph{no dynamics ever connects them});
\item[(iv)] $\delta\chi_S(z)=0$ for every $z\notin\operatorname{spec}M$
\hfill (\emph{exact zero at all frequencies}).
\end{enumerate}
A system satisfying these conditions is in \textbf{Class I}; we set
$\nu=\infty$.
\end{theorem}

\subsection{Why it is true (sketch)}

\emph{(i)$\Rightarrow$(iii).}  Cayley--Hamilton (Tool 2): once the first $n$
moments vanish, all higher powers of $M$ are combinations of lower ones, so
all moments vanish, hence $\ell^\dagger q(M)r=0$ for any polynomial.

\emph{(iii)$\Rightarrow$(iv).}  The resolvent $(zI-M)^{-1}$ is itself a
polynomial in $M$ (with $z$-dependent coefficients): by Cramer's rule
$(zI-M)^{-1}=\adj(zI-M)/\det(zI-M)$, and the adjugate of $zI-M$ expands as a
polynomial in $M$.  So (iii) kills $\delta\chi_S(z)$ at every $z$ off the
spectrum --- \emph{not} only at large $|z|$.

\emph{(iv)$\Rightarrow$(i).}  If a Laurent series
$\sum_k\mu_k z^{-k-1}$ vanishes on an open annulus, all its coefficients
vanish (uniqueness of Laurent expansions).

\emph{(ii)$\Leftrightarrow$(iii)}: unwinding definitions --- $\kry(M,r)$ is
spanned by the vectors $M^k r$, and membership in $\nobs$ says exactly that
$\ell^\dagger M^j$ kills them for every $j$.

\begin{takeaway}
An exact all-frequency zero is not a miracle of fine-tuning at each
frequency; it is a \emph{structural} statement --- the reachable subspace is
invisible to the detector --- and it can be \emph{certified by computing just
$n$ numbers}.  This finite certificate is Step~1 of the decision algorithm.
\end{takeaway}

\subsection{Four mechanisms, one condition}

Physically distinct causes all funnel into condition (ii):
\begin{enumerate}
\item \textbf{Symmetry separation}: source and detector live in different
symmetry sectors that the dynamics never mixes (selection rules);
\item \textbf{Unreachability}: the input never spreads into the part of the
space the detector watches;
\item \textbf{Unobservability}: the input spreads everywhere, but the
detector is blind to the whole invariant subspace it spreads in;
\item \textbf{All-order cancellation}: no invariant subspace is visible in
the physical basis, yet the $n$ moments cancel ``accidentally.''
\end{enumerate}
The theorem does not care which mechanism operates --- and that is a feature:
the certificate is the same finite computation in all four cases.

\begin{example}[A $2\times2$ Class I system]
Take
\[
M=\begin{pmatrix}0&1\\0&0\end{pmatrix},\qquad
r=\begin{pmatrix}1\\0\end{pmatrix},\qquad
\ell=\begin{pmatrix}0\\1\end{pmatrix}.
\]
Then $Mr=0$, so the reachable subspace is
$\kry(M,r)=\spann\{r\}=\spann\{e_1\}$, and $\ell^\dagger e_1=0$: everything
reachable is invisible.  Check the moments: $\mu_0=\ell^\dagger r=0$,
$\mu_1=\ell^\dagger Mr=0$ --- both vanish, and $n=2$, so Theorem~\ref{thm:one}
promises an exact zero.  Indeed
\[
(zI-M)^{-1}=\begin{pmatrix}1/z&1/z^2\\0&1/z\end{pmatrix}
\quad\Longrightarrow\quad
\ell^\dagger(zI-M)^{-1}r=0\ \text{ for every }z\neq0 .
\]
Note the drama hidden in this tiny example: $M$ \emph{does} couple the two
levels ($M_{12}=1\neq0$), but in the wrong direction --- level 2 feeds
level 1, never the reverse.  Connectivity of the matrix is not enough; what
counts is directed reachability from $r$ to $\ell$.
\end{example}

% ==================================================================
\section{Theorem II: the dissipative hierarchy (Class II, $0<\nu<\infty$)}
\label{sec:thm2}
% ==================================================================

\subsection{Setting: invertible dissipation}

Now assume the dissipation matrix $D$ is \emph{invertible} on the relevant
space: every mode feels the damping when we crank $\Gamma$ up.  (The opposite
case --- some modes feel nothing --- is Theorem~III.)  Intuitively,
everything should die as $\Gamma\to\infty$.  Theorem~II makes that precise
and, more importantly, tells you \emph{how fast}, with an integer exponent
that has a direct physical meaning.

\subsection{The expansion}

Factor the large parameter out of $A_\Gamma=\Gamma D+B(z)$:
\[
A_\Gamma = D\,\big(\Gamma I + D^{-1}B(z)\big)
\quad\Longrightarrow\quad
A_\Gamma^{-1}
=\frac1\Gamma\Big(I+\frac{D^{-1}B(z)}{\Gamma}\Big)^{-1} D^{-1}.
\]
For $\Gamma$ larger than $X_K:=\sup_{z\in K}\|D^{-1}B(z)\|$, Tool 1 applies
with $Y=D^{-1}B(z)/\Gamma$ and gives a convergent series in $1/\Gamma$:

\begin{theorem}[Dissipative asymptotic hierarchy]
\label{thm:two}
Let $D$ be invertible and define the \emph{dissipative moments}
\[
\mu_k(z)=p^\dagger\big[D^{-1}B(z)\big]^k D^{-1}c,\qquad k\ge0 .
\]
Then for all $\Gamma>X_K$ and all $z$ in the compact window $K$,
\begin{equation}
F_\Gamma(z)=\sum_{k=0}^{\infty}(-1)^k\,\Gamma^{-(k+1)}\,\mu_k(z),
\label{eq:hierarchy}
\end{equation}
absolutely and uniformly convergent on $K$, with an explicit geometric bound
on every truncation error.  If the first $m$ moments vanish identically on
$K$ and $\mu_m\not\equiv0$, then
\[
F_\Gamma(z)=(-1)^m\,\mu_m(z)\,\Gamma^{-(m+1)}+O\!\big(\Gamma^{-(m+2)}\big)
\qquad\text{uniformly on }K,
\]
and the \emph{dissipative suppression index}
\[
\nudiss:=m+1
\]
is well defined, basis independent, and bounded by the dimension $N$ of the
response-relevant space: $\nudiss\in\{1,\dots,N\}$.  If instead \emph{all}
moments vanish, then $F_\Gamma\equiv0$ and the system is actually in
Class~I.  Systems with $0<\nudiss<\infty$ form \textbf{Class II}, with
$\nu=\nudiss$.
\end{theorem}

Three remarks on why each claim holds:
\begin{itemize}
\item \emph{Convergence and the error bound} are verbatim the geometric
series: the $k$-th term is bounded by $\mathrm{const}\cdot(X_K/\Gamma)^k/\Gamma$.
\item \emph{``$m$ cannot exceed $N-1$''} is Cayley--Hamilton again (Tool 2),
applied to the matrix $D^{-1}B(z)$ at a fixed frequency: if the first $N$
moments vanish, all do.  So either a finite $m\le N-1$ exists, or the
response is identically zero.  There is no in-between, no
``vanishes to all orders but is still nonzero'' --- that possibility is
mathematically excluded in finite dimension.
\item \emph{Basis independence}: $F_\Gamma$ itself is unchanged by any
similarity transformation applied consistently to $(D,B,c,p)$, and the
coefficients of an asymptotic expansion of a fixed function are unique;
hence $m$, and with it $\nudiss$, cannot depend on the basis.
\end{itemize}

\subsection{The meaning of the exponent: counting doors}

Why should $\nu$ be an \emph{integer}, and what does it count?  Suppose the
space splits into sectors (source sector, intermediate sectors, readout
sector) such that $D$ and the diagonal part of $B$ preserve each sector,
while the off-diagonal part $V$ of $B$ hops between adjacent sectors.  Each
term of moment $\mu_k$ is a word with $k$ letters $D^{-1}B$; a word needs at
least $d$ hop-letters to carry the source sector into the readout sector,
where $d$ is the graph distance between them.  Hence
\[
\mu_k\equiv0\ \text{ for } k<d
\qquad\Longrightarrow\qquad \nudiss\ge d+1 .
\]
\begin{takeaway}
In Class II, the decay exponent is a \emph{path length}: the number of damped
sectors the signal must traverse, after all symmetry cancellations at lower
order are accounted for.  Every ``door'' costs one factor of $1/\Gamma$.
This is why the old binary criterion was too coarse: it lumped
``suppressed like $1/\Gamma$'' and ``suppressed like $1/\Gamma^{5}$''
into the same verdict, although they are wildly different experimentally.
\end{takeaway}

\begin{example}[A $2\times2$ Class II system, solved exactly]
\label{ex:classII}
Take $N=2$,
\[
D=I,\qquad
B(z)=\begin{pmatrix}z & g\\ g & z\end{pmatrix},\qquad
c=\begin{pmatrix}1\\0\end{pmatrix},\qquad
p=\begin{pmatrix}0\\1\end{pmatrix},
\]
with coupling $g\neq0$: inject in level 1, read out level 2.  The moments:
\[
\mu_0=p^\dagger D^{-1}c=p^\dagger c=0,
\qquad
\mu_1(z)=p^\dagger B(z)\,c = g\neq0 .
\]
So $m=1$ and Theorem~\ref{thm:two} predicts $\nu=\nudiss=2$.  We can verify
this without any expansion, because the model is exactly solvable:
\[
(\Gamma I+B)^{-1}
=\frac{1}{(\Gamma+z)^2-g^2}
\begin{pmatrix}\Gamma+z&-g\\-g&\Gamma+z\end{pmatrix}
\ \Longrightarrow\
F_\Gamma(z)=\frac{-g}{(\Gamma+z)^2-g^2}
\;\xrightarrow[\Gamma\to\infty]{}\;
-\frac{g}{\Gamma^{2}}+O(\Gamma^{-3}).
\]
Exactly the predicted $\Gamma^{-2}$, with the predicted coefficient
$(-1)^1\mu_1=-g$.  The exponent 2 is ``one hop ($d=1$) plus one'': the signal
enters the damped level 1 (one factor $1/\Gamma$) and must hop once to reach
level 2 (a second factor).
\end{example}

% ==================================================================
\section{Theorem III: protected channels (Class III, $\nu=0$)}
\label{sec:thm3}
% ==================================================================

\subsection{Setting: singular dissipation and dark states}

Theorem~II assumed every mode feels the damping.  But the most interesting
quantum-optical systems violate exactly this: they possess \emph{dark
states} --- superpositions that the dissipation cannot touch.  In matrix
language, $D$ is \emph{singular}: it has a nonzero kernel, $\ker D\neq\{0\}$.
Cranking up $\Gamma$ then does nothing at all to the kernel directions.  The
physical expectation is that a response routed entirely through the kernel
should \emph{survive} $\Gamma\to\infty$.  Theorem~III confirms this and
computes the surviving limit in closed form.

\subsection{Splitting the space: the Riesz projection}

We need to split $\mathcal V$ into ``protected'' (kernel of $D$) and
``damped'' (the rest).  If $D$ were Hermitian we would use the orthogonal
projection onto $\ker D$.  But $D$ lives in Liouville space and is generally
\emph{not} Hermitian, so its eigenspaces need not be orthogonal --- the right
notion is the \emph{spectral} (Riesz) projection: the projection $P$ onto
$\ker D$ \emph{along} the span of all other eigenspaces.
It satisfies $P^2=P$ and $PD=DP=0$, but in general $P^\dagger\ne P$: the
splitting $\mathcal V = P\mathcal V\oplus Q\mathcal V$ (with $Q=I-P$) is
direct but \emph{oblique}, like resolving a vector along two non-perpendicular
axes.  For the curious: $P$ has the elegant contour-integral formula
$P=\frac{1}{2\pi i}\oint_\gamma(\zeta I-D)^{-1}d\zeta$ with $\gamma$ a small
circle around the eigenvalue $0$; for a diagonalizable $D$ you may simply
think ``keep the zero-eigenvalue part of the eigen-decomposition.''

One technical hypothesis is needed: the zero eigenvalue must be
\emph{semisimple} (no Jordan block --- geometric multiplicity equals algebraic
multiplicity), so that $D$ restricted to the damped part, $D_Q$, is honestly
invertible.  Lindblad dissipators in the models of this repository satisfy
this.

\subsection{The theorem}

Write $B_{XY}=XB(z)Y$ for the blocks of $B$ in the oblique splitting, and let
$B_P(z):=PB(z)P|_{P\mathcal V}$ be the \emph{protected block}: the coherent
dynamics that the kernel directions experience among themselves.

\begin{theorem}[Protected channel under singular dissipation]
\label{thm:three}
Assume on the compact window $K$: the zero eigenvalue of $D$ is semisimple,
the protected block $B_P(z)$ is invertible for every $z\in K$, and $B$ is
bounded on $K$.  Then there is $\Gamma_*<\infty$ such that for $\Gamma>\Gamma_*$,
\begin{equation}
F_\Gamma(z)=F_0(z)+\Gamma^{-1}F_1(z)+O(\Gamma^{-2})
\qquad\text{uniformly on }K,
\label{eq:protexp}
\end{equation}
with leading coefficient
\begin{equation}
\boxed{\;F_0(z)=p^\dagger P\, B_P(z)^{-1}\, P c\;}
\label{eq:F0}
\end{equation}
If $F_0\not\equiv0$ on $K$, the response survives arbitrarily strong
dissipation at order one: \textbf{Class III}, $\nu=0$.
\end{theorem}

Formula \eqref{eq:F0} is worth savoring.  It says: in the limit of infinite
noise, the system behaves as if only the dark subspace existed --- project
the source into it ($Pc$), evolve with the protected dynamics alone
($B_P^{-1}$), project onto the detector ($p^\dagger P$).  The infinitely
damped sector has not merely become small; it has become \emph{invisible},
and \eqref{eq:F0} is the transfer function of the effective closed system
living in the kernel.

\subsection{Why it is true (sketch)}

Block-partition $A_\Gamma=\Gamma D+B$ in the splitting
$P\mathcal V\oplus Q\mathcal V$.  Because $PD=DP=0$, the huge term
$\Gamma D$ appears \emph{only} in the damped-damped block:
\[
A_\Gamma=\begin{pmatrix}
B_{PP} & B_{PQ}\\
B_{QP} & \Gamma D_Q+B_{QQ}
\end{pmatrix}.
\]
Now use block elimination (Tool 3): solve the damped block and substitute.
The damped block's inverse is
$(\Gamma D_Q+B_{QQ})^{-1}=\Gamma^{-1}D_Q^{-1}+O(\Gamma^{-2})$ --- \emph{small},
of order $1/\Gamma$.  Hence the Schur complement of the protected block is
\[
S_P(\Gamma,z)=B_{PP}-B_{PQ}\,(\Gamma D_Q+B_{QQ})^{-1}B_{QP}
= B_P(z)+O(\Gamma^{-1}),
\]
and inverting it (legitimate since $B_P$ is invertible and the correction is
small) gives $x_P = B_P^{-1}Pc + O(\Gamma^{-1})$, while the damped component
$x_Q=O(\Gamma^{-1})$.  Contracting with $p^\dagger$ yields
\eqref{eq:protexp}--\eqref{eq:F0}.  The first correction $F_1$ can be written
out explicitly (it contains the leakage terms $B_{PQ},B_{QP}$), and the
expansion can be pushed to any order.

\begin{example}[A $2\times2$ Class III system, solved exactly]
Take
\[
D=\begin{pmatrix}0&0\\0&1\end{pmatrix},\qquad
B(z)=\begin{pmatrix} z & g\\ g & z+i\end{pmatrix},\qquad
c=p=\begin{pmatrix}1\\0\end{pmatrix},
\]
on a window $K$ of frequencies bounded away from the points where $B_P$
degenerates.  Here $\ker D=\spann\{e_1\}$ is the dark direction, $P=\mathrm{diag}(1,0)$,
and the protected block is the $1\times1$ matrix $B_P(z)=z$, invertible for
$z\neq0$.  Formula \eqref{eq:F0} predicts $F_0(z)=1/z$.  Exactly:
\[
F_\Gamma(z)=\big[(\Gamma D+B)^{-1}\big]_{11}
=\frac{\Gamma+z+i}{\,z\,(\Gamma+z+i)-g^2\,}
\;\xrightarrow[\Gamma\to\infty]{}\;\frac1z .
\]
The coupling $g$ to the damped level has \emph{not} been assumed zero ---
leakage out of the dark state is allowed!  It only enters at order
$1/\Gamma$: divide numerator and denominator by $\Gamma$ to see
$F_\Gamma=1/(z-g^2/(\Gamma{+}z{+}i))=1/z+O(\Gamma^{-1})$.  Protection is a
statement about the \emph{leading order}, not about perfect isolation.
\end{example}

\subsection{A trap: touching the dark state is not enough}
\label{sec:trap}

It is tempting to shortcut the theorem: ``the source overlaps the dark
subspace ($Pc\neq0$) and the detector sees it ($p^\dagger P\neq0$), so surely
the response is protected.''  \textbf{False.}  The protected transfer
\eqref{eq:F0} sandwiches $B_P^{-1}$ between the endpoint projections, and
that matrix inverse can have a zero exactly where you sit.  Explicitly, on a
two-dimensional dark subspace take
\[
B_P=\begin{pmatrix}0&1\\1&0\end{pmatrix}
\ \ (\text{so } B_P^{-1}=B_P),\qquad
Pc = \begin{pmatrix}1\\0\end{pmatrix},\qquad
P^\dagger p = \begin{pmatrix}1\\0\end{pmatrix}
\ \Longrightarrow\
F_0 = \begin{pmatrix}1&0\end{pmatrix}B_P\begin{pmatrix}1\\0\end{pmatrix}=0,
\]
even though both endpoint overlaps are nonzero.  A selection rule inside the
dark subspace kills the leading order, and the class is then decided by the
subleading coefficients $F_1,F_2,\dots$ --- the system silently drops into
Class II (or even Class I).  Moral: \emph{the classification cannot be read
off from endpoint overlaps; you must evaluate the protected transfer
function itself}.  This lemma is what makes the decision algorithm of
Section~\ref{sec:algorithm} genuinely necessary.

% ==================================================================
\section{The unified index $\nu$ and the exclusive trichotomy}
\label{sec:trichotomy}
% ==================================================================

So far, three theorems in three regimes.  The unification step defines one
number that summarizes them all.

\begin{definition}[Response-class index]
On the window $K$, define
\[
\nu:=\sup\Big\{\,s\ge0:\ \sup_{z\in K}\Gamma^{s}\,|F_\Gamma(z)|
\xrightarrow[\Gamma\to\infty]{}0\Big\}\ \in[0,\infty],
\]
with $\nu=\infty$ when $F_\Gamma$ vanishes identically at large $\Gamma$.
\end{definition}

In words: $\nu$ is the best power of $\Gamma$ you can multiply the response
by and still see it die.  If the response is exactly zero you can use any
power ($\nu=\infty$); if it decays like $\Gamma^{-m}$ you can use any
$s<m$ but not $s>m$ (so $\nu=m$); if it tends to a nonzero constant even
$s=0^+$ fails ($\nu=0$).

\begin{theorem}[Exclusive and exhaustive trichotomy]
Within the declared scope (Section~\ref{sec:scope}), exactly one of the
following holds:
\[
\begin{array}{llll}
\textbf{Class I} & F_\Gamma(z)\equiv0 & \nu=\infty
& \text{(exact structural no-go)}\\[2pt]
\textbf{Class II} & F_\Gamma(z)=\Gamma^{-m}r_m(z)+O(\Gamma^{-m-1}),\ m\ge1 & \nu=m
& \text{(asymptotic suppression)}\\[2pt]
\textbf{Class III} & F_\Gamma(z)=r_0(z)+O(\Gamma^{-1}),\ r_0\not\equiv0 & \nu=0
& \text{(protected survival)}
\end{array}
\]
\end{theorem}

Why \emph{exclusive}?  Because the leading term of an asymptotic expansion of
a fixed function is unique: if $\Gamma^{m}F_\Gamma(z_0)$ converges to a
nonzero number for one $m$, it diverges or vanishes for every other $m$.
One model, one function $F_\Gamma$, one value of $\nu$ --- a model cannot
carry two class labels.  (And $\nu$ is basis independent, because $F_\Gamma$
itself is: any similarity transformation applied consistently to
$(D,B,c,p)$ cancels inside $p^\dagger A_\Gamma^{-1}c$.)

Why \emph{exhaustive}?  Run the case analysis.  If $D$ is invertible,
Theorem~\ref{thm:two} applies: either all moments vanish (Class I) or a
first nonzero moment exists at some finite order $\le N-1$ (Class II) ---
Cayley--Hamilton forbids anything else.  If $D$ is singular with a semisimple
kernel, Theorem~\ref{thm:three} gives a power series starting at order
$\Gamma^0$: its first nonvanishing coefficient is either the zeroth
(Class III), a later one (Class II), or there is none (Class I).  Every
branch of the tree ends in exactly one of the three classes.

% ==================================================================
\section{The decision algorithm: classification in finitely many steps}
\label{sec:algorithm}
% ==================================================================

Everything above is constructive, and it assembles into an algorithm.  The
punchline --- remarkable for a statement about the behavior of a function at
\emph{all} frequencies and \emph{all} large dissipation strengths --- is that
the algorithm needs only \emph{finitely many exact linear-algebra
operations}.

\medskip
\noindent\textbf{Input:} the matrices $(D,B(z),c,p)$ on the response-relevant
space (dimension $N$), and a realization $(M,r,\ell)$ of the master response
(dimension $n$).

\begin{enumerate}
\item[\textbf{Step 1.}] \emph{Test for Class I.}  Compute the $n$ moments
$\mu_k=\ell^\dagger M^k r$, $k=0,\dots,n-1$.  All zero $\Rightarrow$
\textbf{Class I}, $\nu=\infty$.  Stop.  (Correctness: Theorem~\ref{thm:one}.)
\item[\textbf{Step 2.}] \emph{If $D$ is invertible.}  Compute the dissipative
moments $\mu_k(z)=p^\dagger(D^{-1}B)^kD^{-1}c$ for $k=0,1,\dots$ until the
first $m$ with $\mu_m\not\equiv0$; Cayley--Hamilton guarantees $m\le N-1$
(otherwise Step 1 would have fired).  Return \textbf{Class II} with
$\nu=m+1$.  (Correctness: Theorem~\ref{thm:two}.)
\item[\textbf{Step 3.}] \emph{If $D$ is singular.}  Verify the zero
eigenvalue is semisimple, build the Riesz projection $P$, and evaluate the
protected transfer $F_0$ of \eqref{eq:F0}.  If $F_0\not\equiv0$: return
\textbf{Class III}, $\nu=0$.  If $F_0\equiv0$: compute the subleading
coefficients $F_1,F_2,\dots$; the first nonzero one gives \textbf{Class II}
with the corresponding index; if all vanish, the Class-I certificate of
Step 1 confirms \textbf{Class I}.  (Correctness: Theorem~\ref{thm:three} and
Section~\ref{sec:trap}.)
\end{enumerate}

Why is Step 3 not an infinite search?  Because by Cramer's rule
$F_\Gamma$ is a \emph{rational} function of $\Gamma$ ---
a ratio of two polynomials
$p^\dagger\adj(\Gamma D+B)\,c\,/\det(\Gamma D+B)$
whose degrees are bounded by the dimension.  A rational function whose first
$2N{+}1$ expansion coefficients vanish is identically zero, and the
suppression order can even be read off directly as a difference of two
polynomial degrees in $\Gamma$.  Finite data, finite answer.

\begin{remark}[Why ``exact'' arithmetic is insisted upon]
The theorems certify \emph{identities} --- e.g.\ ``this function is exactly
zero at every frequency.''  Floating-point samples can make an identity
plausible but can never prove it: no finite set of approximate evaluations
certifies a statement about a continuum.  The rigorous documents therefore
require matrix entries in an exact function class (rational or algebraic
expressions with symbolic parameters) where ``is this expression identically
zero?'' is decidable.  Numerics remain invaluable --- the repository verifies
all three theorems numerically --- but as \emph{corroboration}, not
certification.
\end{remark}

% ==================================================================
\section{The fine print: what the theory does \emph{not} say}
\label{sec:scope}
% ==================================================================

An honest no-go theorem lists its rules.  The classification holds within,
and only within, this declared scope:

\begin{itemize}
\item[(A1)] \textbf{Finite dimension.}  The state space is finite
dimensional.  Continua of modes (free-space photon fields, phonon baths kept
explicitly) are excluded.
\item[(A2)] \textbf{Markovian dynamics.}  The environment is memoryless
(Lindblad form).  Non-Markovian memory kernels are excluded.
\item[(A3)] \textbf{Weak probe / linear response.}  Everything is first
order in the probe field.  Strong-probe nonlinearities are excluded.
\item[(A4)] \textbf{No gain.}  The system is passive on the physical
frequency domain, so the resolvents exist where needed.
\item[(A5)] \textbf{A fixed scaling path.}  ``Strong dissipation'' means the
specific one-parameter family $\Gamma\mapsto\Gamma D+B$: one designated
group of dissipators is scaled together, everything else held fixed.
Change what you scale, and you change the question --- and possibly the
class.  The split of the dynamics into $D$ versus $B$ is physical input,
not a mathematical convenience.
\item[(A6)] \textbf{Semisimple kernel (for Theorem III).}  If the zero
eigenvalue of $D$ has a nontrivial Jordan block, the clean power series in
$1/\Gamma$ can fail (fractional powers may appear), and the analysis needs
different tools (Newton polygons).  Such models are declared out of scope.
\end{itemize}

Consequently, the phrase ``material-agnostic'' in the research documents has
a precise, limited meaning: \emph{within} finite-dimensional Markovian
weak-probe response, the classification does not care which material or
platform realizes the matrices.  It does \emph{not} claim to classify every
conceivable open-system model.

\begin{sloppypar}
Also outside the theory's claims: whether a Class III response is
\emph{experimentally visible} (linewidths, optical depth, ensemble
averaging, and noise are all separate questions), and the spectroscopic
distinction between EIT and Autler--Townes splitting.  The rigorous
documents list these explicitly as follow-up physics tasks, distinct from
the closed mathematical classification.
\end{sloppypar}

% ==================================================================
\section{Summary and dictionary}
% ==================================================================

\subsection{One-paragraph summary}

Take an open quantum system in the standard weak-probe setting, designate an
internal sector, and form the master response
$R_{S,\Gamma}(z)=\chi_{\mathrm{full}}-\chi_{\mathrm{cut}}$ --- the part of
the signal that dies when the sector is unplugged.  As the scaled dissipation
$\Gamma$ grows, exactly one of three things happens to this single scalar
function: it is identically zero by structure (Class I, certified by $n$
moment checks via Cayley--Hamilton), it decays like $\Gamma^{-\nu}$ with an
integer $\nu$ counting the damped sectors the signal must traverse
(Class II, via the geometric series in $1/\Gamma$), or it survives at order
one on a dissipation-free dark subspace (Class III, via block elimination
around the kernel of $D$).  The index $\nu\in\{\infty\}\cup(0,\infty)\cup\{0\}$
is basis independent, the three classes are mutually exclusive and
exhaustive, and a three-step, finitely-terminating algorithm decides the
class of any given model.

\subsection{Jargon dictionary}

\begin{center}
\begin{tabular}{p{0.32\textwidth}p{0.60\textwidth}}
\textbf{Term} & \textbf{Plain meaning}\\[3pt]
\hline\\[-8pt]
Liouville space & Density matrices flattened into column vectors, so
superoperators become ordinary matrices.\\[3pt]
Transfer function & The complex number turning ``how hard you poke'' into
``how much you hear,'' as a function of frequency.\\[3pt]
Sector cut & Surgically deleting the couplings of one internal pathway
(inside $B$ only) to see what the detector loses.\\[3pt]
Master response $R_{S,\Gamma}$ & The full-minus-cut difference; the one
object the whole theory classifies.\\[3pt]
Realization $(M,r,\ell)$ & A state-space presentation
$\ell^\dagger(zI-M)^{-1}r$ of a transfer function.\\[3pt]
Moments $\mu_k$ & The numbers $\ell^\dagger M^k r$: overlap between
detector and $k$-times-evolved source.\\[3pt]
Reachable subspace $\kry(M,r)$ & Everywhere the injected excitation can
spread under the dynamics.\\[3pt]
Unobservable subspace\newline $\nobs(M,\ell)$ & Everything the detector can
never see, no matter how long it waits.\\[3pt]
Dark state / kernel of $D$ & A direction the scaled dissipation cannot
touch: $Dv=0$.\\[3pt]
Riesz projection $P$ & The (generally non-orthogonal) projection onto the
dark subspace along the damped one.\\[3pt]
Protected block $B_P$ & The coherent dynamics the dark directions
experience among themselves; its inverse is the surviving response.\\[3pt]
Schur complement & The effective matrix for the slow variables after the
fast ones are eliminated.\\[3pt]
Index $\nu$ & The decay exponent of the response at strong dissipation:
$\infty$, a positive integer, or $0$.\\[3pt]
\hline
\end{tabular}
\end{center}

\subsection{Where to go next}

\begin{itemize}
\item The rigorous proofs of Theorems I--III, the trichotomy, and the
algorithm: the file \path{three_theorems_proofs.tex} in this folder.
\item The closed ``theorem package'' formulation, in which everything is
stated for the master response in its doubled realization, with the
representation-invariance theorem and the polynomial certificate:
the companion master-response document.
\item Numerical verification of all three theorems and of the
endpoint-overlap counterexample of Section~\ref{sec:trap}:
\path{src/three_theorem_verification.py} in this repository
(four PASS lines: a machine-precision Class-I zero, a fitted Class-II slope
$-\nudiss$, a Class-III plateau at $F_0$ with $1/\Gamma$ correction, and the
$F_0=0$ counterexample).
\item Background reading: Kalman controllability/observability (any
linear-systems textbook); the quantum Zeno effect and decoherence-free
subspaces (for physical cousins of Class II and Class III); singular
perturbation theory (for the mathematical cousin of Theorem III).
\end{itemize}

\end{document}
